\section[菱形上的 étale 上同调]{菱形上的 étale 上同调（12 月 14 日）}

\srcpage{72}

我们的目标是为定义 $D(\Bun_G,\Zl)$（$\ell\neq p$）奠定基础。因此需要一套关于一般小 v-叠的 $\ell$-进 étale 层理论。这套理论已在 [S17] 中、在挠系数的情形下发展出来。今天我们将概述它。

\subsection{设定}

回忆 $\Perf$ 是特征 $p$ 的完满空间范畴。我们在其上赋予 v-拓扑。现在有一个技术问题：$\Perf$ 是大的。因此引入小 v-层（small v-sheaf）的概念是有用的；它是可表层的小余极限。等价地，一个 v-层 $\cF$ 是小的，当且仅当存在某个完满空间 $X$ 及满射 $X\to\cF$。

\begin{remark}
利用这样的表示，可以看出小 v-层的子 v-层自动是小的。
\end{remark}

类似地，还有「小 v-叠」（small v-stack）的概念。

\subsection{总体思路}

口号：我们总要用下降化归到严格完全间断空间的情形。

先验地看，这很奇怪：严格完全间断空间相当不几何，例如它们不享有任何 Poincaré 对偶。然而这套办法确实有效。

\subsection{形式化}

设 $\Lambda$ 是某个与 $p$ 互素的整数 $n$ 所湮灭的环。目标：对任意小 v-叠，定义一个三角 $\Lambda$-线性范畴 $D_{\et}(X,\Lambda)$，以及下列 6 个函子：

\begin{enumerate}[label=(\arabic*)]
\item 对任意小 v-叠之间的映射 $f\colon Y\to X$，有拉回
\[
f^*\colon D_{\et}(X,\Lambda)\to D_{\et}(Y,\Lambda),
\]
及其右伴随
\[
f_*\colon D_{\et}(Y,\Lambda)\to D_{\et}(X,\Lambda).
\]
\item 每个 $D_{\et}(X,\Lambda)$ 都有对称幺半张量积 $-\overset{L}{\otimes}_{\Lambda}-$，且 $f^*$ 保持 $-\overset{L}{\otimes}_{\Lambda}-$（即是对称幺半的）。它有部分右伴随
\[
\mathrm{RHom}_{\Lambda}(-,-)\colon D_{\et}(X,\Lambda)^{\mathrm{op}}\times D_{\et}(X,\Lambda)\to D_{\et}(X,\Lambda),
\]
使得
\[
\mathrm{RHom}_{D_{\et}(X,\Lambda)}(A,\mathrm{RHom}_{\Lambda}(B,C))=\mathrm{RHom}_{D_{\et}(X,\Lambda)}(A\overset{L}{\otimes}_{\Lambda}B,C).
\]
\item 若 $f\colon Y\to X$ 在局部空间菱形中可表、可紧化\footnote{在解析几何中，只要紧化存在，就总有典范的紧化，所以这在实践中并不难验证。}，且 $\dim.\mathrm{tr.g.}(f)<\infty$（局部地），则有函子
\[
Rf_!\colon D_{\et}(Y,\Lambda)\to D_{\et}(X,\Lambda),
\]
它满足基变换与投影公式
\[
Rf_!(A\overset{L}{\otimes}_{\Lambda}f^*B)\cong(Rf_!A)\overset{L}{\otimes}_{\Lambda}B
\]
对 $A\in D_{\et}(Y,\Lambda)$、$B\in D_{\et}(X,\Lambda)$ 成立。它容许右伴随
\[
Rf^{!}\colon D_{\et}(X,\Lambda)\to D_{\et}(Y,\Lambda).
\]
\end{enumerate}

\srcpage{73}

实际上我们只需要三个函子 $f^*$、$-\overset{L}{\otimes}_{\Lambda}-$、$Rf_!$；其余函子随后作为右伴随得到。在这三个函子中，我们需要 $f^*$ 与 $-\overset{L}{\otimes}_{\Lambda}-$ 交换，并且与 $Rf_!$ 满足如下意义的基变换：若有笛卡尔图表
\[
\begin{tikzcd}
Y' \arrow[r, "g'"] \arrow[d, "f'"'] & Y \arrow[d, "f"] \\
X' \arrow[r, "g"'] & X
\end{tikzcd}
\]
使得 $Rf_!$ 有定义，则
\[
g^*Rf_!\xrightarrow{\ \sim\ }R(f')_!(g')^*
\]
作为 $D_{\et}(Y,\Lambda)\to D_{\et}(X',\Lambda)$ 的函子同构。

事实上，所有 $D_{\et}(X,\Lambda)$ 都作为稳定 $\Lambda$-线性 ∞-范畴的同伦范畴出现；所有函子都在这个层面上定义。（然而，我们并不声称对 $D_{\et}(X,\Lambda)$ 发展了完整的六函子形式化。这样做有一些技术困难。）

\subsection{导出范畴的定义}

函子 $X\mapsto D_{\et}(X,\Lambda)$ 将是一个 ∞-范畴的 v-层。特别地，若 $Y\to X$ 是 v-覆盖，则
\[
D_{\et}(X,\Lambda)\xrightarrow{\ \sim\ }\varprojlim\left(D_{\et}(Y,\Lambda)\rightrightarrows D_{\et}(Y\times_X Y,\Lambda)\Rrightarrow\cdots\right).
\]

\begin{remark}
函子 $X\mapsto D_{\et}(X,\Lambda)$ 甚至是超 v-层（hyper-v-sheaf），即对超覆盖 $Y_\bullet\to X$ 有类似的陈述。
\end{remark}

正因如此，只需对严格完全间断的完满空间 $X$ 定义 $D_{\et}(X,\Lambda)$。

\begin{definition}
设 $X$ 严格完全间断。定义 $D_{\et}(X,\Lambda):=D(X_{\et},\Lambda)$，即 $X$ 上 étale $\Lambda$-模的 Abel 范畴的导出 ∞-范畴。由一般的 site 论形式论证，它自动继承：
\begin{itemize}
\item 对称幺半张量积 $-\overset{L}{\otimes}_{\Lambda}-$；
\item $*$-拉回函子性。
\end{itemize}
\end{definition}
由此，对任意小 v-叠 $X$，我们已经得到良定义的范畴 $D_{\et}(X,\Lambda)$，以及 $-\overset{L}{\otimes}_{\Lambda}-$ 与 $f^*$。

\begin{proposition}
函子 $f^*$ 有右伴随 $Rf_*$，而 $-\overset{L}{\otimes}_{\Lambda}-$ 有部分右伴随 $\mathrm{RHom}_{\Lambda}(-,-)$。
\end{proposition}

\begin{proof}
对严格完全间断的 $X$，$D_{\et}(X,\Lambda)$ 是可表现（presentable）稳定 ∞-范畴，因为它是一个 site 上的层范畴。于是对一般的 v-叠 $X$，$D_{\et}(X,\Lambda)$ 是可表现 ∞-范畴的小极限，因而也是可表现的。为构造这些右伴随，调用 Lurie 的伴随函子定理。
\end{proof}

现在我们已经定义了 $D_{\et}(X,\Lambda)$，但希望更具体地理解它。设 $X$ 是局部空间菱形。则它有 étale 拓扑 $X_{\et}$，定义如下。映射 $f\colon Y\to X$ 称为 étale 的，若它局部分离，且对一切完满空间 $X'\to X$，映射 $Y\times_X X'\to X'$ 可由一个在 $X'$ 上 étale 的完满空间表出。理想的情形是：对这种 $X$，$D_{\et}(X,\Lambda)$ 就是该拓扑的导出范畴。这几乎是对的。

\srcpage{74}

\begin{warning}
请密切注意下文的上下标。

按定义，$D(X_{\et},\Lambda)$ 是 $X$ 上 étale $\Lambda$-模的 Abel 范畴的导出 ∞-范畴，而 $D_{\et}(X,\Lambda)$ 是上面通过下降定义的范畴。
\end{warning}

\begin{theorem}
存在自然函子
\[
D(X_{\et},\Lambda)\to D_{\et}(X,\Lambda),
\]
它诱导等价
\[
D^+(X_{\et},\Lambda)\to D^+_{\et}(X,\Lambda),
\]
并把 $D_{\et}(X,\Lambda)$ 实现为左完备化
\[
\varprojlim_n D^{\ge -n}_{\et}(X,\Lambda)\cong\varprojlim_n D^{\ge -n}(X_{\et},\Lambda).
\]
特别地，若 $D(X_{\et},\Lambda)$ 左完备，例如在「$X_{\et}$ 的有限上同调维数」下，则
\[
D(X_{\et},\Lambda)\xrightarrow{\ \sim\ }D_{\et}(X,\Lambda).
\]
\end{theorem}

\begin{proof}[证明概要]
第一个观察是：$D_{\et}(X,\Lambda)$ 总有自然的 t-结构（由严格完全间断空间下降而来），且该 t-结构左完备。（关于左完备性，化归到严格完全间断的 $X$；此时上同调维数为 $0$，容易证明。）

因此只需证明 $D^{\ge0}(X_{\et},\Lambda)\xrightarrow{\sim}D^{\ge0}_{\et}(X,\Lambda)$。我们断言后者作为全子范畴包含于 $D^{\ge0}(X_v,\Lambda)$，即 $D^{\ge0}_{\et}(X,\Lambda)\hookrightarrow D^{\ge0}(X_v,\Lambda)$。这是因为 $D^{\ge0}_{\et}(X,\Lambda)$ 是沿 v-超覆盖的极限，而这样的数据按定义下降为 v-拓扑上的数据。
\[
\begin{tikzcd}
D^{\ge0}(X_{\et},\Lambda) \arrow[r] \arrow[dr, dashed] & D^{\ge0}_{\et}(X,\Lambda) \arrow[d, hook] \\
& D^{\ge0}(X_v,\Lambda)
\end{tikzcd}
\]
现在关键是证明：对 $\lambda\colon X_v\to X_{\et}$，$\lambda^*$ 全忠实。等价地，对一切 étale $\Lambda$-层 $\cF$，
\[
\cF\xrightarrow{\ \sim\ }\lambda_*\lambda^*\cF,
\]
且 $R^i\lambda_*\lambda^*\cF=0$（$i>0$）。（口号：「上同调在从 étale 拓扑过渡到 v-拓扑时不变。」）这是 Grothendieck 关于 étale 层的 étale 上同调与 fppf 乃至 fpqc 上同调一致的结果的类比。

为证明这一点，我们先从 étale 拓扑过渡到 pro-étale 拓扑。这基本上是形式的：把 pro-étale 覆盖写成 étale 覆盖的余滤过极限。于是上同调上得到滤过余极限，而滤过余极限是正合的。

接着从 pro-étale 拓扑过渡到 v-拓扑。由 pro-étale 下降，可以假设 $X$ 严格完全间断。若 $Y\to X$ 是由仿射完满空间给出的 v-覆盖，则可写成 $Y=\varprojlim_i Y_i\to X$，其中每个 $Y_i\to X$ 是 $X$ 上某个有限维球中的开集。（这是把任意仿射概形逼近为有限型概形之极限的类比；但这里甚至能把这种逼近安排成由「光滑」对象给出！）若 $X=\Spa(R,R^+)$、$Y=\Spa(S,S^+)$，则存在某个 $I$ 与满射 $R\langle T_i^{1/p^\infty}:i\in I\rangle\twoheadrightarrow S$（例如取 $I$ 为 $S$ 的一切幂有界元素之集）。这把 $S$ 表示成 $S=V(f_j:j\in J)\subset\mathbb{B}^I_X$。它是
\[
\{|f_j|\le|\epsilon|,\ j\in J'\subseteq J\}\subset\mathbb{B}^{I'}_X
\]
\srcpage{75}
沿有限 $I'\subseteq I$ 与有限 $J'\subset J$ 的极限。

于是我们取这些作为 $Y_i$。注意这类似于 Zariski 闭浸入曾是开邻域上的极限。

于是 $Y\to X$ 的下降化归为 $Y_i\to X$ 的下降。（这里用了所考虑的层来自 étale 拓扑，以得知它们在极限上的截面是截面的余极限。）但后者有截面：这些 $Y_i$ 在 $X$ 上是「光滑」的，因此在 $X$ 的某个 étale 覆盖之后有满射；而由于 $X$ 严格完全间断，它没有非平凡的 étale 覆盖。

口号：v-世界中的 Noether 下降容许事物是光滑的！
\end{proof}

\begin{proposition}
对一切小 v-叠 $X$，
\[
D_{\et}(X,\Lambda)\hookrightarrow D(X_v,\Lambda),
\]
且 $A\in D(X_v,\Lambda)$ 落在像中，当且仅当所有 $H^i(A)$ 都落在像中。此外，这可以在拉回到任意局部空间菱形 $Y\to X$ 之后检验；此时「落在像中」的意思是该对象来自一个 étale 层。
\end{proposition}

\subsection{基变换}

\begin{theorem}[qcqs 基变换]
考虑局部空间菱形的笛卡尔图表
\[
\begin{tikzcd}
Y' \arrow[r, "g'"] \arrow[d, "f'"'] & Y \arrow[d, "f"] \\
X' \arrow[r, "g"'] & X
\end{tikzcd}
\]
其中 $f$ 是 qcqs 的。则
\[
Rf'_*(g')^*\xrightarrow{\ \sim\ }g^*Rf_*\colon D^+_{\et}(Y,\Lambda)\to D^+_{\et}(X',\Lambda).
\]
若 $f$ 有有限上同调维数，则这在 $D_{\et}(Y,\Lambda)\to D_{\et}(X',\Lambda)$ 上成立。
\end{theorem}

\begin{proof}
若 $X'\to X$ 是 pro-étale 的，则基变换自动成立。因此可以假设 $X$、$X'$ 严格局部。再利用 $Y$ 上的下降，还可以假设 $Y$ 严格完全间断。于是通过过渡到连通分支，也可假设 $Y$ 严格局部。所以
\[
X=\Spa(C,C^+),\quad X'=\Spa(C',(C')^+),\quad Y=\Spa(\breve C,\breve C^+),\quad Y'=\Spa(R,R^+),
\]
且 $R=C'\widehat{\otimes}_C\breve C\supset R^+$。
\end{proof}

\begin{lemma}
我们有
\[
H^i(Y',\mathbb{Z}/n\mathbb{Z})=
\begin{cases}
\mathbb{Z}/n\mathbb{Z}, & i=0,\\
0, & i>-0.
\end{cases}
\]
\end{lemma}

这可以由上同调在代数闭域扩张下的不变性推出；对此，可以化归到有限型情形下 Huber 的工作。\yizhu{原文印作 $i>-0$，疑为 $i>0$ 之误；此处照录，未改原文。}{A（原文 p.75）}\hfill$\square$

注意这里不需要紧合性。但这不是你习惯的那种基变换，因为 adic 空间的「点」不是 adic 空间。例如对严格局部空间 $\Spa(C,C^+)$，我们并没有沿 $\{s\}\hookrightarrow\Spa(C,C^+)$ 的基变换，因为 $\{s\}$ 不是 adic 空间，甚至不是菱形。

\begin{theorem}[紧合基变换]
若 $f$ 还是紧合的，则我们也得到这类基变换。
\end{theorem}

\begin{proof}[证明概要]
通过类似的化归步骤，这化归为概形的紧合基变换。
\end{proof}

\srcpage{76}

\subsection{紧合前推}

若 $f\colon Y\to X$ 在空间菱形中可表、可紧化，且 $\dim.\mathrm{tr.g.}f<\infty$，则有开浸入
\[
\begin{tikzcd}
Y \arrow[rr, hook, "j"] \arrow[dr] & & \bar Y^{/X} \arrow[dl] \\
& X &
\end{tikzcd}
\]
其中 $\bar Y^{/X}$ 由 $\bar Y^{/X}(R,R^+)=X(R,R^+)\times_{X(R,R^\circ)}Y(R,R^\circ)$ 给出。这可以用赋值判别法来表述，即沿 $\Spa(R,R^\circ)\to\Spa(R,R^+)$ 延拓映射。回忆 $Y\to X$ 紧合当且仅当
\[
\begin{tikzcd}
\Spa(R,R^\circ) \arrow[r] \arrow[d] & Y \arrow[d] \\
\Spa(R,R^+) \arrow[r] \arrow[ur, dashed, "\exists!"] & X
\end{tikzcd}
\]
这等价于 $Y\xrightarrow{\ \sim\ }\bar Y^{/X}$。对上面的一般 $Y\to X$，$\bar Y^{/X}$ 是 $X$ 上带有来自 $Y$ 的映射的始初紧合菱形。

\begin{definition}
令 $Rf_!=R\bar f^{/X}_*\circ j_!\colon D_{\et}(Y,\Lambda)\to D_{\et}(X,\Lambda)$，其中 $j_!$ 是零延拓（即 $j^*$ 的左伴随）。
\end{definition}

\begin{example}
令 $Y=\mathbb{B}_C\to X=\Spa C$。试着把 $\mathbb{B}_C$ 嵌入一个紧化，例如 $\mathbb{P}^1_C$。
\end{example}

\srcpage{77}
\begin{center}
\figP{077}
\end{center}

典范紧化 $\bar{\mathbb{B}}^{/C}_C$ 添入一个半径为无穷小地大于 $1$ 的秩 $2$ 点。另一种看法是它缩小了 $R^+$：我们有 $\mathbb{B}_C=\Spa(C\langle T\rangle,\cO_C\langle T\rangle)$ 与 $\bar{\mathbb{B}}^{/C}_C=\Spa(C\langle T\rangle,\cO_C+\mathfrak{m}_C\langle T\rangle)$。

一般地，这个紧化「添入所有更高秩的点」。它完全是函子性的。更一般地，只要 $f$ 仅在局部空间菱形中可表（并且可紧化、$\dim.\mathrm{tr.g.}(f)<\infty$），也可以定义它。

\begin{theorem}
$Rf_!$ 满足基变换与投影公式，并且可复合。
\end{theorem}

这其实由紧合基变换容易推出；关键在于 $j_!$ 与 $Rf_*$ 的交换。

\subsection{例外原像}

利用伴随函子定理可以验证：

\begin{proposition}
$Rf_!$ 有右伴随 $Rf^!$。
\end{proposition}

\subsection{Verdier 对偶}

这是最棘手的部分。我们不能直接定义光滑性，因为不存在非既约的完满环，因而没有切空间的概念。

\begin{definition}
设 $f\colon Y\to X$ 在局部空间菱形中可表、可紧化，且 $\dim.\mathrm{tr.g.}(f)<\infty$。若对一切 $A\in D_{\et}(X,\Lambda)$ 有 $Rf^!A\cong f^*A\otimes\mathbb{D}_f$，其中 $\mathbb{D}_f\in D_{\et}(Y,\Lambda)$，且这在任意基变换后仍成立，则称 $f$ 是上同调光滑的。\yizhu{原文此处写作 $Rf^!\cong f^*A\otimes\mathbb{D}_f$，漏写了 $Rf^!$ 作用的对象 $A$；这里按 Verdier 对偶的通常陈述补出。}{A（原文 p.77）}
\end{definition}

\srcpage{78}

我们是在把一个关于通常 $\ell$-进层的定理（Verdier 对偶）变成定义。困难之处在于找到例子。

\begin{example}
设 $S$ 是无限基数的射有限集。则映射 $f\colon\underline{S}\times\Spa C=\Spa\operatorname{Cont}(S,C)\to\Spa C$ 是 pro-étale 的，但不是上同调光滑的。$Rf^!\Lambda$ 是什么？按定义，
\[
R\Gamma(S,Rf^!\Lambda)=\mathrm{RHom}(Rf_!\Lambda,\Lambda)=\mathrm{RHom}(Rf_*\Lambda,\Lambda)=\operatorname{Cont}(S,\Lambda)^\vee.
\]
这是 $S$ 上 $\Lambda$-值测度的空间。类似地可以计算茎
\[
(Rf^!\Lambda)_s=\varinjlim_{U\ni s}\mathcal{M}(U,\Lambda),
\]
它是「$s$ 处的测度芽」，而不是 $\Lambda$。
\end{example}

\begin{example}
我们可以把上同调光滑的概念推广到叠：要求光滑局部的上同调光滑性。若 $G$ 局部 pro-$p$，则映射 $[\mathrm{pt}/G]\to *$ 是上同调光滑的。这反映了 $G$ 上存在 Haar 测度。
\end{example}
